\section{The constructions and their finite data}
A graph is dismantled by deleting vertices until one vertex remains. In an
ordinary dismantling, each deleted vertex is dominated by a neighbour. The
recursive hierarchy permits a deletion at level $k$ when the graph induced by
its current neighbours is dismantlable at level $k-1$. A graph has
\emph{strict level $k$} if it admits a reduction at level $k$ and none at level
$k-1$.

The article \emph{Linear-size graphs at every strict dismantling level}
constructs graphs $F_k$ and $H_k$ of strict level $k$, for every $k\geq1$,
with
\[
 |V(F_k)|=3k+8,\qquad |V(H_k)|=4k+4.
\]
The title, authors and current article reference are recorded in
\href{https://github.com/Geometric-and-Combinatorial-Structures/strictly-k-dismantlable-graphs/blob/v1.1.1/ARTICLE.md}{\texttt{ARTICLE.md}}.
This supplement develops the finite proof for $H_k$, gives a third family
$J_k$ of order $6k+6$, and explains the kernel completion and homological
identities behind the layer construction. Let $m_k$ be the least order
of a graph of strict level $k$. The two families give
$m_k\leq\min\{4k+4,3k+8\}$; they do not determine $m_k$ in general.
References here to sections, figures and results of the article concern the
identified 23-page manuscript; \texttt{ARTICLE.md} distinguishes it from
the version at the public article locator.

Each family starts with a finite seed. A new layer contains a centre whose
neighbourhood graph is the previous graph, together with auxiliary vertices
that exclude first deletions at the lower level. To apply this construction,
we need a reduction of the seed at its stated level, an obstruction to the
preceding level, a reduction after deleting the first centre, and obstructions
in the remaining neighbourhoods. We supply the finite data for these tasks
and use the homological argument of Section~\ref{sec:homology} to justify the
criterion for every subsequent layer.

For example, $H_1$ has eight vertices. Its reduction at level one requires
only two outer deletions, each justified by a single fold in a neighbourhood;
the remaining graph is a cone. None of its original neighbourhoods is a cone,
so no level-zero reduction can begin. Section~\ref{sec:family4} gives this
example in full before treating the larger recursive certificates.

\noindent\textbf{Organisation.}
Sections~\ref{sec:prelim} and~\ref{sec:construction} fix the hierarchy and
recall the common construction and its Euler criterion. Section~\ref{sec:family4} supplies
the $4k+4$ family, including its two prefixes; Section~\ref{sec:family6}
checks a second application with five auxiliary types. Section~\ref{sec:kernels}
relates constructions with two auxiliary types to an edge-deleted kernel and
identifies the completion producing $F_1$. Its first subsection supplies the
homological preliminaries for Sections~\ref{sec:homology} and~\ref{sec:minimum}.
Section~\ref{sec:homology} proves
the degree shift in non-central neighbourhoods and the resulting homological
criterion. Section~\ref{sec:minimum} explains the finite verification that the minimum at
strict level two is twelve. Appendix~\ref{sec:article} records the finite calculations and
figures for the article's $3k+8$ family.

\noindent\textbf{Conventions and notation.}
The layer indices form $\mathcal L_m$ (written $L_m$ in the article); clique counts are
written $\kappa_t$, reserving $c_i$ for centres. The symbols $F_k$ and $H_k$
retain their meanings in the article; $J_k$ denotes the $6k+6$ family.
Within a section, $A$ denotes its seed, $B$ its first extension and $G_m$ the
graph with $m$ added layers; this layer count is distinct from the minimum
order $m_k$. The seed vertex $a_i$ has numerical label $i$ in
the data files. Homology is
reduced homology over $\mathbb F_2$, with augmented chains in degree $-1$.
The command following each finite calculation runs from the root of the
accompanying repository and requires Python 3.11 or later.
